\chapter{عبارات باقاعده (Regular Expressions)}

در فصول پیش‌رو، به بررسی ابزارهایی خواهیم پرداخت که منحصراً برای پردازش متن طراحی شده‌اند. همان‌طور که پیش‌تر مشاهده کردیم، داده‌های متنی در تمامی سیستم‌های شبه‌یونیکس و لینوکس نقشی حیاتی ایفا می‌کنند. با این حال، پیش از آنکه بتوانیم از تمامی ظرفیت‌های این ابزارها بهره‌مند شویم، باید با فناوری زیربنایی آن‌ها یعنی «عبارات باقاعده» (\en{Regular Expressions}) آشنا شویم؛ فناوری‌ای که اغلب در پیچیده‌ترین سناریوهای پردازش متن به کار گرفته می‌شود.

در مسیر یادگیری قابلیت‌های مختلف خط فرمان، با مفاهیم و دستورات چالش‌برانگیزی نظیر بسط پوسته (\en{shell expansion})، استفاده از کوتیشن (\en{quoting})، میان‌برهای صفحه‌کلید، تاریخچه دستورات و ویرایشگر \cmd{vi} روبه‌رو شده‌ایم. عبارات باقاعده این «سنت» دشواری را ادامه می‌دهند و بی‌شک می‌توان گفت مبهم‌ترین ویژگی در میان آن‌ها هستند. البته این به معنای اتلاف وقت در یادگیری آن‌ها نیست؛ برعکس، تسلط بر عبارات باقاعده ما را قادر به انجام عملیات قدرتمند و پیشرفته‌ای در پردازش داده می‌کند که ارزش واقعی آن‌ها در پروژه‌های بزرگ نمایان می‌شود.

\section{اصول و چیستی عبارات باقاعده (Regular Expressions)}

به زبان ساده، عبارات باقاعده سیستم‌های نشانه‌گذاری نمادینی هستند که برای شناسایی و استخراج الگوهای متنی (\en{Pattern Matching}) به کار می‌روند. این عبارات تا حدودی به نویسه‌های جانشین (\en{wildcards}) در پوسته شباهت دارند که برای تطبیق نام فایل‌ها استفاده می‌کردیم، اما از نظر گستردگی و قدرت عملیاتی در مقیاسی بسیار وسیع‌تر عمل می‌کنند. اکثر ابزارهای خط فرمان و زبان‌های برنامه‌نویسی برای حل چالش‌های پردازش متن از این فناوری پشتیبانی می‌کنند. با این حال، برای درک دقیق‌تر باید توجه داشت که همه‌ی عبارات باقاعده یکسان نیستند و بسته به ابزار یا زبان برنامه‌نویسی، تفاوت‌های جزئی در نگارش آن‌ها وجود دارد. در این مبحث، تمرکز ما بر عبارات باقاعده مطابق با استاندارد \en{POSIX} است (که اکثر ابزارهای خط فرمان لینوکس را پوشش می‌دهد)، در حالی که بسیاری از زبان‌های برنامه‌نویسی (به‌ویژه \en{Perl}) از مجموعه‌ی گسترده‌تر و غنی‌تری از نشانه‌گذاری‌ها استفاده می‌کنند.

\section{جست‌وجوی متنی با ابزار grep}

ابزار محوری ما برای کار با عبارات باقاعده، برنامه نام‌آشنای \cmd{grep} است. عنوان «\en{grep}» در واقع کوته‌نوشتی از عبارت «\en{global regular expression print}» است که نشان‌دهنده پیوند عمیق و ساختاری این ابزار با مفاهیم عبارات باقاعده است. در واقعیت، \cmd{grep} فایل‌های متنی را به منظور یافتن الگوهایی که با یک عبارت باقاعده مشخص تطبیق دارند، پیمایش کرده و هر خطی را که حاوی آن تطبیق (\en{Match}) باشد، در خروجی استاندارد چاپ می‌کند.

تاکنون، ما از \cmd{grep} صرفاً برای جست‌وجوی رشته‌های ثابت (\en{Fixed Strings}) به نحو زیر بهره برده‌ایم:

\begin{latin}
\begin{lstlisting}
[me@linuxbox ~]$ ls /usr/bin | grep zip
\end{lstlisting}
\end{latin}

این دستور، تمامی فایل‌های موجود در دایرکتوری \file{/usr/bin} را که نام آن‌ها حاوی زیررشته‌ی «\en{zip}» باشد، فهرست می‌کند.

ساختار کلی فراخوانی برنامه \cmd{grep} به صورت زیر است، که در آن \file{regex} بیانگر یک عبارت باقاعده (\en{Regular Expression}) می‌باشد:

\begin{latin}
\begin{lstlisting}
grep [options] regex [file...]
\end{lstlisting}
\end{latin}

جدول زیر گزینه‌های پرکاربرد \cmd{grep} را توصیف می‌کند:

\begin{table}[htbp]
\centering
\caption{گزینه‌های پرکاربرد دستور \cmd{grep}}
\label{tab:grep-options}
\begin{tabularx}{\textwidth}{l l X}
\toprule
\textbf{گزینه کوتاه} & \textbf{گزینه بلند} & \textbf{شرح عملکرد} \\
\midrule
\cmd{-i} & \cmd{--ignore-case} & نادیده‌گرفتن حساسیت به حروف: میان حروف کوچک و بزرگ در فرآیند تطبیق تفاوت قائل نمی‌شود. \\
\addlinespace
\cmd{-v} & \cmd{--invert-match} & معکوس‌سازی تطبیق: به‌جای چاپ خطوط دارای تطبیق، تمامی خطوطی را که فاقد الگوی موردنظر هستند نمایش می‌دهد. \\
\addlinespace
\cmd{-c} & \cmd{--count} & شمارش تعداد: به‌جای استخراج خطوط، تعداد کل موارد منطبق (یا نامنطبق در صورت ترکیب با \cmd{-v}) را چاپ می‌کند. \\
\addlinespace
\cmd{-l} & \cmd{--files-with-matches} & فهرست فایل‌های منطبق: به‌جای نمایش محتوا، تنها نام فایل‌هایی را که حداقل یک مورد تطبیق دارند چاپ می‌کند. \\
\addlinespace
\cmd{-L} & \cmd{--files-without-match} & فهرست فایل‌های نامنطبق: مشابه گزینه \cmd{-l} عمل می‌کند، اما تنها نام فایل‌هایی را چاپ می‌کند که فاقد الگوی موردنظر هستند. \\
\addlinespace
\cmd{-n} & \cmd{--line-number} & نمایش شماره خط: در خروجی، پیش از هر خط منطبق، شماره ردیف آن در فایل اصلی را درج می‌کند. \\
\addlinespace
\cmd{-h} & \cmd{--no-filename} & عدم نمایش نام فایل: در جست‌وجوی هم‌زمان میان چندین فایل، از درج نام فایل در ابتدای هر خط خروجی جلوگیری می‌کند. \\
\addlinespace
\cmd{-q} & \cmd{--quiet, --silent} & تمامی خروجی‌های معمول را پنهان می‌کند. این گزینه در اسکریپت‌نویسی پوسته برای بررسی وجود یا عدم وجود یک الگو در فایل‌ها، بدون نیاز به نمایش خروجی، به کار می‌رود؛ به‌گونه‌ای که خروجی دستور صرفاً از طریق کد وضعیت خروجی (\en{Exit Status}) قابل‌تشخیص است. \\
\bottomrule
\end{tabularx}
\end{table}

به منظور بررسی عمیق‌تر عملکرد \cmd{grep} در سناریوهای عملیاتی، ابتدا مجموعه‌ای از فایل‌های متنی حاوی فهرست دایرکتوری‌های سیستمی ایجاد می‌کنیم:

\begin{latin}
\begin{lstlisting}
[me@linuxbox ~]$ ls /bin > dirlist-bin.txt
[me@linuxbox ~]$ ls /usr/bin > dirlist-usr-bin.txt
[me@linuxbox ~]$ ls /sbin > dirlist-sbin.txt
[me@linuxbox ~]$ ls /usr/sbin > dirlist-usr-sbin.txt
[me@linuxbox ~]$ ls dirlist*.txt
dirlist-bin.txt  dirlist-sbin.txt  dirlist-usr-sbin.txt  dirlist-usr-bin.txt
\end{lstlisting}
\end{latin}

اکنون می‌توانیم یک جست‌وجوی متنی ساده را در میان این فایل‌ها اجرا کنیم:

\begin{latin}
\begin{lstlisting}
[me@linuxbox ~]$ grep bzip dirlist*.txt
dirlist-bin.txt:bzip2
dirlist-bin.txt:bzip2recover
\end{lstlisting}
\end{latin}

در این مثال، ابزار \cmd{grep} تمامی فایل‌های هدف را برای یافتن رشته‌ی \en{``bzip''} پیمایش کرده و دو مورد تطبیق یافته را که هر دو در فایل \file{dirlist-bin.txt} مستقر هستند، گزارش می‌کند. چنانچه صرفاً به دریافت فهرست فایل‌های حاوی الگوی مورد نظر (بدون جزئیات خطوط) نیاز داشته باشیم، از گزینه‌ی \cmd{-l} استفاده می‌کنیم:

\begin{latin}
\begin{lstlisting}
[me@linuxbox ~]$ grep -l bzip dirlist*.txt
dirlist-bin.txt
\end{lstlisting}
\end{latin}

در مقابل، اگر هدف شناسایی فایل‌هایی باشد که فاقد الگوی مورد نظر هستند، گزینه‌ی \cmd{-L} به کار گرفته می‌شود:

\begin{latin}
\begin{lstlisting}
[me@linuxbox ~]$ grep -L bzip dirlist*.txt
dirlist-sbin.txt
dirlist-usr-bin.txt
dirlist-usr-sbin.txt
\end{lstlisting}
\end{latin}

\section{متاکاراکترها و نویسه‌های تحت‌اللفظی (Literal Characters)}

شاید در نگاه اول بدیهی به نظر نرسد، اما تمامی جست‌وجوهای ما با ابزار \cmd{grep} تاکنون مبتنی بر عبارات باقاعده بوده‌اند؛ هرچند در ساده‌ترین شکل ممکن. عبارت باقاعده \cmd{bzip} به این معناست که تطبیق تنها زمانی رخ می‌دهد که در یک خط، نویسه‌های \cmd{z}، \cmd{i}، \cmd{p} و \cmd{b} دقیقاً با همین ترتیب و بدون هیچ فاصله‌ای میان آن‌ها ظاهر شوند. نویسه‌های موجود در رشته‌ی \cmd{bzip} همگی «نویسه‌های تحت‌اللفظی» (\en{Literal Characters}) هستند؛ یعنی دقیقاً با خودشان تطبیق داده می‌شوند.

در مقابل، عبارات باقاعده می‌توانند حاوی «متاکاراکترها» (\en{Metacharacters}) باشند که برای تعریف الگوهای پیچیده‌تر و هوشمند به کار می‌روند. این متاکاراکترها عبارتند از:

\begin{latin}
\begin{lstlisting}
^ $ . [ ] { } - ? * + ( ) | \
\end{lstlisting}
\end{latin}

تمامی نویسه‌های دیگر «تحت‌اللفظی» تلقی می‌شوند، مگر نویسه‌ی «بک‌اسلش» (\cmd{\textbackslash}) که در چند مورد خاص برای ایجاد توالی‌های متا (\en{meta sequences}) و همچنین برای گریز دادن متاکاراکترها و رفتار آن‌ها به عنوان نویسه‌های تحت‌اللفظی (به جای تفسیر به عنوان متاکاراکتر) استفاده می‌شود.

\begin{note}
\textbf{نکته:} همان‌طور که مشهود است، بسیاری از متاکاراکترهای عبارات باقاعده با نویسه‌هایی که در «بسط پوسته» (\en{Shell Expansion}) معنای خاص دارند، مشترک هستند. بنابراین، هنگام وارد کردن عبارات باقاعده حاوی متاکاراکتر در ترمینال، بسیار حیاتی است که کل عبارت را درون کوتیشن (\cmd{' '}) قرار دهید تا از تفسیر زودهنگام و نادرست آن توسط پوسته جلوگیری شود.
\end{note}

\section{نویسه نقطه (The Any Character)}

نخستین متاکاراکتری که بررسی می‌کنیم، نقطه (\cmd{.}) یا همان نویسه‌ی \en{Period} است. این نماد برای تطبیق با «هر نویسه‌ی دلخواهی» به کار می‌رود. چنانچه این متاکاراکتر را در یک عبارت باقاعده قرار دهید، هر نویسه‌ای که در آن موقعیتِ کاراکتری قرار گرفته باشد، موجب برقراری تطبیق (\en{Match}) خواهد شد. به مثال زیر توجه کنید:

\begin{latin}
\begin{lstlisting}
[me@linuxbox ~]$ grep -h '.zip' dirlist*.txt
bunzip2
bzip2
bzip2recover
gunzip
gzip
funzip
gpg-zip
preunzip
prezip
prezip-bin
unzip
unzipsfx
\end{lstlisting}
\end{latin}

در این نمونه، ما به دنبال تمامی خطوطی بودیم که با الگوی \cmd{.zip} مطابقت داشته باشند. خروجی حاصل، حاوی نکات فنی قابل‌تأملی است؛ نخست اینکه برنامه \cmd{zip} در فهرست نتایج مشاهده نمی‌شود. علت این امر آن است که وجود متاکاراکتر نقطه در ابتدای عبارت باقاعده، طول رشته‌ی مورد انتظار برای تطبیق را به حداقل چهار نویسه افزایش داده است، به طوری که حتماً باید یک نویسه (هر چه که باشد) پیش از حرف \cmd{z} قرار داشته باشد.

علاوه بر این، اگر در فایل‌های مورد جست‌وجو، نامی با پسوند \file{.zip} وجود داشت، آن نیز در خروجی ظاهر می‌شد؛ چرا که در منطق عبارات باقاعده، متاکاراکتر نقطه می‌تواند خودِ نویسه‌ی نقطه (نقطه‌ی واقعی در نام فایل) را نیز به عنوان «یک نویسه‌ی دلخواه» پوشش داده و تطبیق دهد.

\section{لنگرها (Anchors)}

در ادبیات عبارات باقاعده، نویسه‌های «کَرِت» (\en{Caret}) که با نماد \cmd{\^{}} نمایش داده می‌شود و «علامت دلار» (\cmd{\$}) به عنوان لنگر (\en{Anchor}) شناخته می‌شوند. وظیفهٔ لنگرها این است که تطبیق الگو را به موقعیت‌های مکانی خاصی در یک خط محدود کنند؛ به این معنا که تطبیق تنها زمانی احراز می‌شود که الگوی مورد نظر در ابتدای خط (\cmd{\^{}}) یا انتهای خط (\cmd{\$}) قرار گرفته باشد.

به مثال‌های زیر دقت کنید:

\begin{latin}
\begin{lstlisting}
[me@linuxbox ~]$ grep -h '^zip' dirlist*.txt
zip
zipcloak
zipgrep
zipinfo
zipnote
zipsplit
[me@linuxbox ~]$ grep -h 'zip$' dirlist*.txt
gunzip
gzip
funzip
gpg-zip
preunzip
prezip
unzip
zip
[me@linuxbox ~]$ grep -h '^zip$' dirlist*.txt
zip
\end{lstlisting}
\end{latin}

در این نمونه‌ها، فهرست فایل‌ها را برای یافتن رشته‌ی \cmd{zip} در سه وضعیت متفاوت بررسی کردیم: ابتدا در آغاز خط، سپس در پایان خط، و در نهایت خطی که \cmd{zip} هم‌زمان هم در ابتدا و هم در انتهای آن ظاهر شده است؛ به عبارت دیگر، خطی که محتوای آن منحصراً کلمه‌ی \cmd{zip} باشد.

\begin{note}
\textbf{نکته:} عبارت باقاعده \cmd{\^{}\$} برای شناسایی خطوط خالی (\en{Empty Line}) به کار می‌رود. این الگو به این معناست که از ابتدای خط (\cmd{\^{}}) مستقیماً به انتهای خط (\cmd{\$}) می‌رسیم، بدون آنکه هیچ نویسه‌ای (حتی فاصله یا تب) در میان باشد. به بیان ساده‌تر، \cmd{\^{}\$} خطوطی را تطبیق می‌دهد که کاملاً خالی هستند و هیچ‌گونه کاراکتر قابل‌مشاهده‌ای در آن‌ها وجود ندارد.
\end{note}

\section{کاربرد عبارات باقاعده در حل جداول کلمات متقاطع}

حتی با تکیه بر دانش مختصری که تا این مرحله از عبارات باقاعده کسب کرده‌ایم، می‌توانیم مسائل کاربردی جالبی را حل کنیم. برای مثال، سناریویی را در نظر بگیرید که در آن نیاز دارید واژه‌ای پنج‌حرفی را بیابید که حرف سوم آن «\en{j}»، حرف آخر آن «\en{r}» و معنای خاصی داشته باشد.

آیا می‌دانستید که سیستم‌عامل لینوکس شما به‌طور پیش‌فرض دارای یک لغت‌نامه جامع است؟ اگر به دایرکتوری \file{/usr/share/dict/} مراجعه کنید، یک یا چند فایل لغت‌نامه (\en{Dictionary}) خواهید یافت. این فایل‌ها در واقع فهرست‌های متنی طویلی هستند که در هر سطر یک واژه را به ترتیب حروف الفبا در خود جای داده‌اند. در بسیاری از توزیع‌ها، فایلی به نام \file{words} وجود دارد که شامل ده‌ها هزار واژه است. برای یافتن پاسخ احتمالی سوال بالا، می‌توان از دستور زیر استفاده کرد:

\begin{latin}
\begin{lstlisting}
[me@linuxbox ~]$ grep -i '^..j.r$' /usr/share/dict/words
Major
major
\end{lstlisting}
\end{latin}

در این الگو، با استفاده از لنگرها (\cmd{\^{}} و \cmd{\$}) و نقطه (\cmd{.}) به عنوان جایگزین حروف مجهول، تمامی واژگان پنج‌حرفی که در موقعیت‌های تعیین‌شده دارای حروف «\en{j}» و «\en{r}» هستند، استخراج شده‌اند.

\section{عبارات براکتی و کلاس‌های نویسه‌ای (Bracket Expressions)}

علاوه بر امکان تطبیق «هر نویسه‌ای» در یک موقعیت خاص، می‌توان با استفاده از «عبارات براکتی»، دامنه تطبیق را تنها به مجموعه‌ی مشخصی از نویسه‌ها محدود کرد. با این روش، می‌توان فهرستی از کاراکترهای مجاز را (حتی متاکاراکترهایی که درون براکت معنای خاص خود را از دست می‌دهند) تعیین نمود. به عنوان نمونه، برای یافتن خطوطی که شامل هر یک از رشته‌های \cmd{bzip} یا \cmd{gzip} هستند، از الگوی زیر استفاده می‌کنیم:

\begin{latin}
\begin{lstlisting}
[me@linuxbox ~]$ grep -h '[bg]zip' dirlist*.txt
bzip2
bzip2recover
gzip
\end{lstlisting}
\end{latin}

یک مجموعه می‌تواند شامل هر تعداد نویسه باشد. نکته مهم این است که اکثر متاکاراکترها هنگام قرارگیری درون براکت، ویژگی نمادین خود را از دست داده و به نویسه‌های تحت‌اللفظی تبدیل می‌شوند. با این حال، دو استثنای حیاتی وجود دارد که در آن‌ها نویسه‌ها همچنان معنای عملیاتی خاصی دارند:

\begin{enumerate}
  \item \textbf{نویسه کَرِت (\cmd{\^{}}):} اگر به عنوان اولین نویسه درون براکت قرار گیرد، برای «نفی» (\en{Negation}) مجموعه به کار می‌رود.
  \item \textbf{خط تیره (\cmd{-}):} برای تعیین یک «محدوده نویسه‌ای» (\en{Character Range}) استفاده می‌شود.
\end{enumerate}

\subsection{نفی در عبارات براکتی (Negation)}

چنانچه نویسه‌ی کَرِت (\cmd{\^{}}) به عنوان نخستین کاراکتر درون یک عبارت براکتی قرار گیرد، مجموعه‌ی نویسه‌های پس از آن به عنوان کاراکترهای «غیرمجاز» در نظر گرفته می‌شوند. به عبارتی، تطبیق تنها زمانی احراز می‌گردد که نویسه‌ی موجود در آن موقعیت، عضوی از مجموعه‌ی نفی‌شده نباشد. با بازنگری در مثال قبلی، این سازوکار به شکل زیر عمل می‌کند:

\begin{latin}
\begin{lstlisting}
[me@linuxbox ~]$ grep -h '[^bg]zip' dirlist*.txt
bunzip2
gunzip
funzip
gpg-zip
preunzip
prezip
prezip-bin
unzip
unzipsfx
\end{lstlisting}
\end{latin}

با اعمال منطق نفی، فهرستی از خطوط استخراج شده است که همگی شامل رشته‌ی \cmd{zip} هستند، اما نویسه‌ی بلافاصله قبل از آن‌ها نباید \cmd{b} یا \cmd{g} باشد. نکته‌ی فنی حائز اهمیت این است که واژه‌ی \cmd{zip} به تنهایی در خروجی ظاهر نشده است؛ زیرا الگو همچنان بر وجود «یک نویسه» پیش از \cmd{zip} تأکید دارد، مشروط بر اینکه آن نویسه عضو مجموعه‌ی نفی‌شده نباشد. در واقع، مجموعه‌ی نفی‌شده جایگاه کاراکتری را حذف نمی‌کند، بلکه محدوده‌ی مقادیر قابل قبول برای آن جایگاه را تغییر می‌دهد.

ذکر این نکته ضروری است که نویسه‌ی \cmd{\^{}} تنها زمانی نقش عملیاتی نفی را ایفا می‌کند که در جایگاه نخستین کاراکترِ درون براکت قرار گیرد. در صورتی که در هر موقعیت دیگری درون براکت درج شود، معنای ویژه‌ی خود را از دست داده و به عنوان یک نویسه‌ی تحت‌اللفظی (\en{Literal}) در مجموعه لحاظ می‌گردد.

\subsection{بازه‌های نویسه‌ای سنتی (Traditional Character Ranges)}

چنانچه بخواهیم یک عبارت باقاعده برای شناسایی فایل‌هایی که با حروف بزرگ آغاز می‌شوند بنویسیم، می‌توانیم تمامی حروف الفبا را مستقیماً درون یک عبارت براکتی درج کنیم:

\begin{latin}
\begin{lstlisting}
[me@linuxbox ~]$ grep -h '^[ABCDEFGHIJKLMNOPQRSTUVWXZY]' dirlist*.txt
\end{lstlisting}
\end{latin}

درج دستی تمامی ۲۶ حرف الفبا، فرآیندی زمان‌بر و مستعد خطا است؛ از این رو، عبارات باقاعده قابلیت تعریف «بازه» یا \en{Range} را فراهم کرده‌اند. با استفاده از این ویژگی، می‌توان دستور فوق را به شکلی فشرده‌تر و کارآمدتر بازنویسی کرد:

\begin{latin}
\begin{lstlisting}
[me@linuxbox ~]$ grep -h '^[A-Z]' dirlist*.txt
MAKEDev
ControlPanel
GET
HEAD
POST
X
X11
Xorg
MAKEFLOPPIES
NetworkManager
NetworkManagerDispatcher
\end{lstlisting}
\end{latin}

در این ساختار، با استفاده از یک الگوی سه‌نویسه‌ای، کل حروف بزرگ الفبا پوشش داده شده است. به همین ترتیب، می‌توان چندین بازه را برای تعریف مجموعه‌های گسترده‌تر ترکیب کرد. برای مثال، الگوی زیر تمامی فایل‌هایی را که نام آن‌ها با یک حرف (بزرگ یا کوچک) یا یک عدد شروع می‌شود، تطبیق می‌دهد:

\begin{latin}
\begin{lstlisting}
[me@linuxbox ~]$ grep -h '^[A-Za-z0-9]' dirlist*.txt
\end{lstlisting}
\end{latin}

در فرآیند تعریف بازه‌ها، نویسه‌ی خط تیره (\cmd{-}) دارای معنای عملیاتی خاصی است. حال این پرسش مطرح می‌شود که اگر بخواهیم خودِ نویسه‌ی خط تیره را به عنوان یکی از اعضای مجموعه جست‌وجو کنیم، چه باید کرد؟ راهکار فنی این است که خط تیره را در «اولین جایگاه» درون براکت قرار دهیم. تفاوت این دو مثال را در نظر بگیرید:

مثال اول، تمامی نام‌فایل‌هایی را که حداقل شامل یک حرف بزرگ هستند، می‌یابد:

\begin{latin}
\begin{lstlisting}
[me@linuxbox ~]$ grep -h '[A-Z]' dirlist*.txt
\end{lstlisting}
\end{latin}

اما در مثال دوم، با قرار دادن خط تیره در ابتدای مجموعه، \cmd{grep} به دنبال خطوطی می‌گردد که حداقل شامل یکی از سه نویسه‌ی «خط تیره»، «\en{A}» یا «\en{Z}» باشند:

\begin{latin}
\begin{lstlisting}
[me@linuxbox ~]$ grep -h '[-AZ]' dirlist*.txt
\end{lstlisting}
\end{latin}

\subsection{کلاس‌های نویسه‌ای POSIX}

بازه‌های نویسه‌ای سنتی (\en{Traditional Character Ranges}) راهکاری بصری و کارآمد برای تعریف سریع مجموعه‌ای از کاراکترها محسوب می‌شوند؛ با این حال، این روش همیشه نتایج قابل‌اتکایی ارائه نمی‌دهد. هرچند تا این مرحله در استفاده از ابزار \cmd{grep} با چالش خاصی مواجه نشده‌ایم، اما ممکن است هنگام کار با سایر برنامه‌ها با رفتارهای غیرمنتظره‌ای روبرو شویم.

در فصل چهارم، نحوه استفاده از نویسه‌های جانشین (\en{Wildcards}) را برای بسط مسیر فایل‌ها بررسی کردیم. در آن مبحث اشاره شد که بازه‌های نویسه‌ای تقریباً مشابه با عبارات باقاعده عمل می‌کنند، اما چالش اصلی در اینجا نهفته است:

\begin{latin}
\begin{lstlisting}
[me@linuxbox ~]$ ls /usr/sbin/[ABCDEFGHIJKLMNOPQRSTUVWXYZ]*
/usr/sbin/MAKEFLOPPIES
/usr/sbin/NetworkManagerDispatcher
/usr/sbin/NetworkManager
\end{lstlisting}
\end{latin}

(بسته به توزیع لینوکس مورد استفاده، ممکن است فهرست متفاوتی دریافت کنید؛ این مثال بر اساس اوبونتو تنظیم شده است). دستور فوق نتیجه‌ی مورد انتظار را تولید می‌کند: فهرستی از فایل‌هایی که دقیقاً با یکی از حروف بزرگ الفبا آغاز می‌شوند. اما با اجرای دستور زیر، نتیجه‌ای کاملاً متفاوت حاصل می‌شود:

\begin{latin}
\begin{lstlisting}
[me@linuxbox ~]$ ls /usr/sbin/[A-Z]*
/usr/sbin/biosdecode
/usr/sbin/chat
/usr/sbin/chgpasswd
/usr/sbin/chpasswd
/usr/sbin/chroot
/usr/sbin/cleanup-info
/usr/sbin/complain
/usr/sbin/console-kit-daemon
\end{lstlisting}
\end{latin}

علت این رخداد چیست؟ پاسخ در پیشینه‌ی تاریخی سیستم‌عامل یونیکس نهفته است. در ابتدای توسعه‌ی یونیکس، تنها استاندارد \en{ASCII} به رسمیت شناخته می‌شد. در این استاندارد، ۳۲ کد نخست (۰ تا ۳۱) به نویسه‌های کنترلی (مانند \en{Tab} و \en{Backspace}) اختصاص داشت. ۳۲ کد بعدی (۳۲ تا ۶۳) شامل نویسه‌های چاپی، علائم نگارشی و ارقام ۰ تا ۹ بود. بازه‌ی بعدی (۶۴ تا ۹۵) حروف بزرگ و برخی علائم دیگر را در بر می‌گرفت و در نهایت، بازه‌ی پایانی (۹۶ تا ۱۲۶) شامل حروف کوچک و باقی‌مانده‌ی علائم نگارشی بود.

بر اساس این ساختار، ترتیب مرتب‌سازی (\en{Collating Sequence}) در سیستم‌های مبتنی بر \en{ASCII} به این صورت تعریف شده بود:

\begin{latin}
\begin{lstlisting}
ABCDEFGHIJKLMNOPQRSTUVWXYZabcdefghijklmnopqrstuvwxyz
\end{lstlisting}
\end{latin}

این ترتیب با منطق رایج در فرهنگ‌های لغت (\en{Dictionary Order}) که حروف را به صورت جفت (کوچک و بزرگ) در کنار هم قرار می‌دهند، تفاوت اساسی داشت:

\begin{latin}
\begin{lstlisting}
aAbBcCdDeEfFgGhHiIjJkKlLmMnNoOpPqQrRsStTuUvVwWxXyYzZ
\end{lstlisting}
\end{latin}

با گسترش جهانی یونیکس و ضرورت پشتیبانی از زبان‌های غیرانگلیسی، جدول \en{ASCII} به ۸ بیت (۲۵۶ کاراکتر) ارتقا یافت. برای مدیریت این تنوع، استانداردهای \en{POSIX} مفهومی تحت عنوان \en{Locale} (تنظیمات بومی) را معرفی کردند تا مجموعه‌ی نویسه‌ها و قوانین مرتب‌سازی بر اساس موقعیت جغرافیایی و زبان کاربر تنظیم شود. شما می‌توانید تنظیمات فعلی سیستم خود را با دستور زیر مشاهده کنید:

\begin{latin}
\begin{lstlisting}
[me@linuxbox ~]$ echo $LANG
en_US.UTF-8
\end{lstlisting}
\end{latin}

در این پیکربندی، برنامه‌های منطبق با استاندارد \en{POSIX} از ترتیب مرتب‌سازی دیکشنری به‌جای ترتیب عددی \en{ASCII} استفاده می‌کنند. این دقیقاً همان دلیلی است که بازه‌ی \cmd{[A-Z]} در مثال قبلی رفتاری غیرمنتظره داشت؛ در ترتیب دیکشنری، این بازه شامل تمامی حروف الفبا (بزرگ و کوچک) به جز حرف کوچک \en{a} می‌شود.

استاندارد \en{POSIX} برای غلبه بر مشکلات ناشی از تفاوت در تنظیمات بومی (\en{Locales})، مجموعه‌ای از کلاس‌های نویسه‌ای استاندارد را معرفی کرده است. این کلاس‌ها انتزاعی از مجموعه کاراکترهای پرکاربرد هستند که فارغ از نوع ترتیب مرتب‌سازی سیستم، همواره خروجی یکسانی ارائه می‌دهند.

\begin{table}[htbp]
\centering
\caption{کلاس‌های نویسه‌ای \en{POSIX}}
\label{tab:posix-classes}
\begin{tabularx}{\textwidth}{l X l}
\toprule
\textbf{کلاس نویسه‌ای} & \textbf{شرح عملکرد فنی} & \textbf{معادل تقریبی در \en{ASCII}} \\
\midrule
\cmd{[:alnum:]} & نویسه‌های الفبایی-عددی (\en{Alphanumeric}) & \cmd{[A-Za-z0-9]} \\
\addlinespace
\cmd{[:word:]} & نویسه‌های واژگانی (شامل حروف، اعداد و زیرخط) — افزونه‌ی \en{GNU} (غیراستاندارد). این کلاس نویسه‌ای در استاندارد \en{POSIX} تعریف نشده و مختص ابزارهای \en{GNU} است. & \cmd{[A-Za-z0-9\_]} \\
\addlinespace
\cmd{[:alpha:]} & نویسه‌های الفبایی (\en{Alphabetic}) & \cmd{[A-Za-z]} \\
\addlinespace
\cmd{[:blank:]} & نویسه‌های تهی (شامل فاصله و \en{Tab}) & \cmd{[ \textbackslash t]} \\
\addlinespace
\cmd{[:cntrl:]} & نویسه‌های کنترلی (غیرقابل چاپ) & کدهای ۰ تا ۳۱ و ۱۲۷ \\
\addlinespace
\cmd{[:digit:]} & ارقام ده‌دهی & \cmd{[0-9]} \\
\addlinespace
\cmd{[:graph:]} & نویسه‌های مرئی (همهٔ نویسه‌های قابل چاپ به جز فاصله) & کدهای ۳۳ تا ۱۲۶ \\
\addlinespace
\cmd{[:lower:]} & حروف کوچک الفبا & \cmd{[a-z]} \\
\addlinespace
\cmd{[:punct:]} & نویسه‌های نشانه‌گذاری و نمادها & \cmd{[-!"\#\$\%\&'()*+,./:;<=>?@[\textbackslash\textbackslash]\_\`\{|\}\textasciitilde]} \\
\addlinespace
\cmd{[:print:]} & نویسه‌های قابل چاپ (شامل نویسه‌های مرئی و فاصله) & کدهای ۳۲ تا ۱۲۶ \\
\addlinespace
\cmd{[:space:]} & فضاهای خالی (\en{Whitespace})؛ شامل فاصله، تب، بازگشت به ابتدای سطر (\en{Carriage Return})، خط جدید، تب عمودی و \en{Form Feed}. & \cmd{[ \textbackslash t\textbackslash r\textbackslash n\textbackslash v\textbackslash f]} \\
\addlinespace
\cmd{[:upper:]} & حروف بزرگ الفبا & \cmd{[A-Z]} \\
\addlinespace
\cmd{[:xdigit:]} & اعداد هگزادسیمال: نویسه‌های مورد استفاده در مبنای ۱۶ (ارقام ۰ تا ۹ و حروف \en{A} تا \en{F}). & \cmd{[0-9A-Fa-f]} \\
\bottomrule
\end{tabularx}
\end{table}

استفاده از این کلاس‌ها دقت فنی و قابلیت حمل (\en{Portability}) اسکریپت‌های شما را تضمین می‌کند. برای مثال، جهت یافتن فایل‌هایی که با حروف بزرگ شروع می‌شوند، به جای \cmd{[A-Z]} از ساختار زیر استفاده می‌کنیم (توجه داشته باشید که کلاس نویسه‌ای خود باید درون یک جفت براکت دیگر قرار گیرد):

\begin{latin}
\begin{lstlisting}
[me@linuxbox ~]$ ls /usr/sbin/[[:upper:]]*
/usr/sbin/MAKEFLOPPIES
/usr/sbin/NetworkManagerDispatcher
/usr/sbin/NetworkManager
\end{lstlisting}
\end{latin}

شایان ذکر است که علی‌رغم مزایای غیرقابل‌انکار این کلاس‌ها، هنوز راهکار استاندارد و منعطفی برای بیان بازه‌های جزئی (مانند بازه‌ی حروف \en{A} تا \en{M}) از طریق این کلاس‌ها وجود ندارد و در چنین موارد خاصی، همچنان باید به محدوده‌های سنتی یا سایر تکنیک‌های تطبیق تکیه کرد.

\section{بازگشت به ترتیب مرتب‌سازی سنتی ASCII}

شما می‌توانید پیکربندی سیستم خود را به‌گونه‌ای تغییر دهید که از ترتیب مرتب‌سازی سنتی (\en{ASCII}) پیروی کند. این امر با اصلاح مقدار متغیر محیطی \cmd{LANG} میسر است. همان‌طور که پیش‌تر اشاره شد، متغیر \cmd{LANG} دربرگیرنده نام زبان و مجموعه‌ی نویسه‌های مورد استفاده در تنظیمات بومی (\en{Locale}) شماست که معمولاً در فرآیند نصب سیستم‌عامل تعیین می‌گردد.

برای مشاهده‌ی جزئیات دقیق تنظیمات بومی فعلی، از دستور \cmd{locale} استفاده کنید:

\begin{latin}
\begin{lstlisting}
[me@linuxbox ~]$ locale
LANG=en_US.UTF-8
LC_CTYPE="en_US.UTF-8"
LC_NUMERIC="en_US.UTF-8"
LC_TIME="en_US.UTF-8"
LC_COLLATE="en_US.UTF-8"
...
LC_ALL=
\end{lstlisting}
\end{latin}

جهت بازگرداندن رفتار سیستم به حالت سنتی یونیکس (پیروی از اولویت‌های عددی \en{ASCII})، باید مقدار متغیر \cmd{LANG} را بر روی \cmd{POSIX} تنظیم نمایید:

\begin{latin}
\begin{lstlisting}
[me@linuxbox ~]$ export LANG=POSIX
\end{lstlisting}
\end{latin}

شایان ذکر است که این تغییر، مجموعه‌ی نویسه‌های سیستم را به انگلیسی آمریکایی (\en{ASCII}) محدود می‌کند؛ لذا پیش از اعمال آن، از انطباق این شرایط با نیازهای خود اطمینان حاصل کنید. در صورتی که تمایل دارید این تغییر به صورت دائمی اعمال شود، می‌توانید عبارت زیر را به فایل پیکربندی \file{.bashrc} در دایرکتوری خانگی خود اضافه کنید:

\begin{latin}
\begin{lstlisting}
export LANG=POSIX
\end{lstlisting}
\end{latin}

با اعمال این تنظیم، بازه‌های نویسه‌ای مانند \cmd{[A-Z]} دقیقاً مطابق با جدول \en{ASCII} تفسیر شده و تنها حروف بزرگ را شامل می‌شوند.

\section{عبارات باقاعده پایه (BRE) و توسعه‌یافته (ERE) در POSIX}

درست در زمانی که به نظر می‌رسد مفاهیم به ثبات رسیده‌اند، با تفکیک پیاده‌سازی عبارات باقاعده توسط استاندارد \en{POSIX} مواجه می‌شویم. این استاندارد، عبارات باقاعده را به دو دسته‌ی اصلی تقسیم می‌کند: عبارات باقاعده پایه (\en{Basic Regular Expressions - BRE}) و عبارات باقاعده توسعه‌یافته (\en{Extended Regular Expressions - ERE}). تمامی ویژگی‌هایی که تا این بخش بررسی کردیم، در برنامه‌های منطبق با استاندارد \en{BRE} (از جمله نسخه‌ی استاندارد \cmd{grep}) پشتیبانی می‌شوند.

تمایز بنیادین میان \en{BRE} و \en{ERE} در نحوه‌ی تفسیر متاکاراکترها نهفته است. در استاندارد \en{BRE}، متاکاراکترهای اصلی و پیش‌فرض عبارتند از:

\begin{latin}
\begin{lstlisting}
^ $ . [ ] *
\end{lstlisting}
\end{latin}

با این حال، نویسه‌های \cmd{( ) \{ \} ? + |} نیز می‌توانند با استفاده از بک‌اسلش (\cmd{\textbackslash}) به متاکاراکتر تبدیل شوند. به عبارت دیگر، در \en{BRE} برای استفاده از قابلیت‌های گروه‌بندی (\cmd{\textbackslash(} و \cmd{\textbackslash)}) و تکرارکننده‌های بازه‌ای (\cmd{\textbackslash\{} و \cmd{\textbackslash\}})، باید از بک‌اسلش استفاده کرد تا این نویسه‌ها نقش ویژه‌ی خود را ایفا کنند. در غیر این صورت، به‌عنوان نویسه‌های تحت‌اللفظی (\en{Literal}) تفسیر می‌شوند. برای نمونه، الگوی \cmd{a\{1,2\}} در \en{BRE} به‌معنای تطابق تحت‌اللفظی با رشته‌ی \cmd{a\{1,2\}} است، در حالی که \cmd{a\textbackslash\{1,2\textbackslash\}} به‌معنای تطابق با \cmd{a} یا \cmd{aa} خواهد بود.

اما در استاندارد \en{ERE}، متاکاراکترهای عملیاتی زیر نیز به فهرست افزوده می‌شوند و بدون نیاز به بک‌اسلش نقش ویژه‌ی خود را ایفا می‌کنند:

\begin{latin}
\begin{lstlisting}
( ) { } ? + |
\end{lstlisting}
\end{latin}

نکته‌ی فنی و چالش‌برانگیز اینجاست: در استاندارد \en{BRE}، نویسه‌های پرانتز \cmd{( )} و آکولاد \cmd{\{ \}} تنها زمانی به عنوان متاکاراکتر عمل می‌کنند که با یک بک‌اسلش (\cmd{\textbackslash}) همراه شوند؛ در حالی که در \en{ERE}، این نویسه‌ها به‌صورت پیش‌فرض متاکاراکتر هستند و افزودن بک‌اسلش به آن‌ها، ویژگی نمادین‌شان را لغو کرده و آن‌ها را به نویسه‌ی تحت‌اللفظی تبدیل می‌کند.

از آنجایی که قابلیت‌هایی که در ادامه‌ی این مبحث بررسی خواهیم کرد متعلق به خانواده‌ی \en{ERE} هستند، باید از نسخه‌ی متفاوتی از \cmd{grep} استفاده کنیم. در سیستم‌های سنتی یونیکس، این وظیفه بر عهده‌ی برنامه‌ی \cmd{egrep} بود؛ اما در اکوسیستم \en{GNU}، ابزار \cmd{grep} با استفاده از گزینه عملیاتی \cmd{-E} به سادگی از عبارات باقاعده توسعه‌یافته پشتیبانی می‌کند.

\section{روند استانداردسازی POSIX}

در دهه‌ی ۱۹۸۰ میلادی، یونیکس به چنان محبوبیتی در دنیای سیستم‌عامل‌های تجاری دست یافت که تا سال ۱۹۸۸، این زیست‌بوم را با آشفتگی و بحرانی جدی مواجه کرد. در آن مقطع، تولیدکنندگان متعدد سخت‌افزار با اخذ مجوز کد منبع یونیکس از آزمایشگاه‌های \en{AT\&T}، نسخه‌های اختصاصی خود را عرضه می‌کردند. با این حال، هر تولیدکننده برای ایجاد مزیت رقابتی و تمایز در بازار، تغییرات و قابلیت‌های انحصاری خود را به سیستم‌عامل می‌افزود. این رویکرد که هدفش وابسته کردن مشتری به پلتفرمی خاص (\en{Vendor Lock-in}) بود، منجر به کاهش شدید سازگاری نرم‌افزاری شد. این دوران سیاه در تاریخ یونیکس، امروزه با اصطلاح «بالکانیزه شدن» (\en{Balkanization}) شناخته می‌شود.

در پاسخ به این ناهماهنگی‌ها، مؤسسه‌ی مهندسان برق و الکترونیک (\en{IEEE}) وارد میدان شد. از اواسط دهه‌ی ۱۹۸۰، این نهاد تدوین مجموعه‌ای از استانداردهای جامع را آغاز کرد تا چارچوب عملکردی یونیکس و سیستم‌های شبه‌یونیکس (\en{Unix-like}) را به دقت ترسیم کند. این استانداردها که به‌طور رسمی با کد \en{IEEE 1003} شناخته می‌شوند، وظیفه‌ی تعریف رابط‌های برنامه‌نویسی (\en{API})، پوسته‌ها (\en{Shells}) و ابزارهای سیستمی را بر عهده دارند که وجود آن‌ها در یک سیستم استاندارد الزامی است.

عنوان \en{POSIX} که کوته‌نوشت عبارت \en{Portable Operating System Interface} است (و حرف \en{X} در انتهای آن صرفاً برای هم‌آوایی با یونیکس و جذابیت بصری اضافه شده)، توسط ریچارد استالمن پیشنهاد و در نهایت از سوی \en{IEEE} تصویب شد. این استاندارد سنگ‌بنای سازگاری و قابلیت انتقال نرم‌افزارها در دنیای لینوکس و متن‌باز امروزی محسوب می‌شود.

\section{تناوب و انتخاب چندگانه (Alternation)}

نخستین ویژگی از «عبارات باقاعده توسعه‌یافته» (\en{ERE}) که بررسی می‌کنیم، قابلیت تناوب است. این ویژگی امکان تطبیق با یکی از چندین عبارت مشخص شده را فراهم می‌آورد. همان‌طور که عبارات براکتی اجازه می‌دهند «یک نویسه» را از میان مجموعه‌ای از کاراکترها انتخاب کنید، تناوب به شما اجازه می‌دهد از میان چندین «رشته» یا «عبارت باقاعده»، یکی را برای تطبیق برگزینید.

برای درک بهتر، از ترکیب دستور \cmd{echo} و \cmd{grep} استفاده می‌کنیم. ابتدا یک جست‌وجوی ساده‌ی رشته‌ای را در نظر بگیرید:

\begin{latin}
\begin{lstlisting}
[me@linuxbox ~]$ echo "AAA" | grep AAA
AAA
[me@linuxbox ~]$ echo "BBB" | grep AAA
[me@linuxbox ~]$
\end{lstlisting}
\end{latin}

در این مثال، خروجی \cmd{echo} به ورودی \cmd{grep} هدایت می‌شود؛ اگر تطبیقی رخ دهد، رشته چاپ شده و در غیر این صورت خروجی خالی می‌ماند. اکنون با استفاده از متاکاراکتر «خط عمودی» (\cmd{|})، قابلیت تناوب را اضافه می‌کنیم:

\begin{latin}
\begin{lstlisting}
[me@linuxbox ~]$ echo "AAA" | grep -E 'AAA|BBB'
AAA
[me@linuxbox ~]$ echo "BBB" | grep -E 'AAA|BBB'
BBB
[me@linuxbox ~]$ echo "CCC" | grep -E 'AAA|BBB'
[me@linuxbox ~]$
\end{lstlisting}
\end{latin}

عبارت باقاعده \cmd{'AAA|BBB'} به معنای «تطبیق با رشته \cmd{AAA} یا رشته \cmd{BBB}» است. توجه داشته باشید که به دلیل استفاده از ویژگی‌های توسعه‌یافته، گزینه \cmd{-E} به دستور \cmd{grep} اضافه شده است. همچنین عبارت باقاعده درون کوتیشن قرار گرفته تا از تفسیر اشتباه کاراکتر \cmd{|} توسط پوسته (به عنوان عملگر \en{Pipe}) جلوگیری شود. تناوب به دو انتخاب محدود نیست و می‌تواند موارد متعددی را در بر بگیرد:

\begin{latin}
\begin{lstlisting}
[me@linuxbox ~]$ echo "AAA" | grep -E 'AAA|BBB|CCC'
AAA
\end{lstlisting}
\end{latin}

برای ترکیب تناوب با سایر مؤلفه‌های عبارات باقاعده، از پرانتز \cmd{()} جهت گروه‌بندی و جداسازی دامنه تناوب استفاده می‌کنیم:

\begin{latin}
\begin{lstlisting}
[me@linuxbox ~]$ grep -Eh '^(bz|gz|zip)' dirlist*.txt
\end{lstlisting}
\end{latin}

این الگو تمامی فایل‌هایی را که دقیقاً با یکی از رشته‌های \cmd{bz}، \cmd{gz} یا \cmd{zip} آغاز می‌شوند، شناسایی می‌کند. اهمیت گروه‌بندی در اینجا بسیار حیاتی است؛ زیرا اگر پرانتزها حذف شوند، منطق عبارت کاملاً تغییر کرده و لنگر ابتدای خط (\cmd{\^{}}) تنها بر روی رشته اول (\cmd{bz}) اعمال می‌شود. در آن صورت، الگو به دنبال هر فایلی می‌گردد که یا با \cmd{bz} شروع شود و یا حاوی رشته‌های \cmd{gz} یا \cmd{zip} در هر کجای نام خود باشد.

\begin{latin}
\begin{lstlisting}
[me@linuxbox ~]$ grep -Eh '^bz|gz|zip' dirlist*.txt
\end{lstlisting}
\end{latin}

\section{کمّی‌کننده‌ها (Quantifiers)}

عبارات باقاعده توسعه‌یافته (\en{ERE}) از چندین روش برای تعیین تعداد دفعات تکرار یک عنصر پشتیبانی می‌کنند. این نمادها به ما اجازه می‌دهند الگوهای انعطاف‌پذیرتری را برای جست‌وجو تعریف کنیم.

\subsection{نماد ? (تطبیق اختیاری صفر یا یک بار)}

این کمّی‌کننده در عمل به معنای «اختیاری کردن عنصر پیشین» است. فرض کنید قصد داریم یک شماره تلفن را اعتبارسنجی کنیم و الگو را زمانی معتبر بدانیم که با یکی از دو فرمت زیر (که در آن‌ها \cmd{n} نشان‌دهنده یک رقم است) مطابقت داشته باشد:

\begin{table}[htbp]
\centering
\caption{ساختارهای معتبر شماره تلفن}
\label{tab:phone-formats}
\begin{tabularx}{\textwidth}{l l}
\toprule
\textbf{ساختار شماره تلفن معتبر (مدل اول)} & \textbf{ساختار شماره تلفن معتبر (مدل دوم)} \\
\midrule
\cmd{(nnn) nnn-nnnn} & \cmd{nnn nnn-nnnn} \\
\bottomrule
\end{tabularx}
\end{table}

برای پوشش هر دو حالت، می‌توانیم عبارتی باقاعده به شکل زیر طراحی کنیم:

\begin{latin}
\begin{lstlisting}
^\(?[0-9][0-9][0-9]\)? [0-9][0-9][0-9]-[0-9][0-9][0-9][0-9]$
\end{lstlisting}
\end{latin}

در این الگو، پس از کاراکترهای پرانتز، علامت سؤال (\cmd{?}) قرار گرفته است تا حضور آن‌ها را اختیاری کند (صفر یا یک بار تکرار). از آنجایی که در استانداردهای توسعه‌یافته، پرانتزها خود متاکاراکتر هستند، از بک‌اسلش (\cmd{\textbackslash}) استفاده شده است تا به عنوان نویسه‌ی تحت‌اللفظی (\en{Literal}) تفسیر شوند.

بیایید الگو را با ابزار \cmd{grep} بیازماییم:

\begin{latin}
\begin{lstlisting}
[me@linuxbox ~]$ echo "(555) 123-4567" | grep -E '^\(?[0-9][0-9][0-9]\)? [0-9][0-9][0-9]-[0-9][0-9][0-9][0-9]$'
(555) 123-4567
[me@linuxbox ~]$ echo "555 123-4567" | grep -E '^\(?[0-9][0-9][0-9]\)? [0-9][0-9][0-9]-[0-9][0-9][0-9][0-9]$'
555 123-4567
[me@linuxbox ~]$ echo "AAA 123-4567" | grep -E '^\(?[0-9][0-9][0-9]\)? [0-9][0-9][0-9]-[0-9][0-9][0-9][0-9]$'
[me@linuxbox ~]$
\end{lstlisting}
\end{latin}

مشاهده می‌کنیم که عبارت طراحی‌شده، هر دو فرمت استاندارد شماره تلفن را شناسایی کرده اما از پذیرش ورودی‌های حاوی کاراکترهای غیرعددی خودداری می‌کند. هرچند این الگو هنوز کامل نیست (زیرا به طور هم‌زمان عدم تقارن در پرانتزها را کنترل نمی‌کند)، اما گام نخست و مؤثری در فرآیند اعتبارسنجی داده محسوب می‌شود.

\subsection{نماد * (تطبیق صفر یا چندباره)}

نماد ستاره (\cmd{*}) مشابه علامت سؤال برای تعریف موارد اختیاری به کار می‌رود؛ با این تفاوت که برخلاف علامت سؤال، عنصر پیشین می‌تواند به هر تعداد دفعات (از صفر تا بی‌نهایت) تکرار شود.

فرض کنید قصد داریم یک ساختار ابتدایی برای تشخیص «جمله» تعریف کنیم؛ به این صورت که رشته با یک حرف بزرگ آغاز شود، سپس با ترکیبی از حروف بزرگ، کوچک و فاصله ادامه یابد و در نهایت به یک نقطه ختم شود. برای پیاده‌سازی این الگو، از عبارت باقاعده زیر استفاده می‌کنیم:

\begin{latin}
\begin{lstlisting}
^[[:upper:]][[:upper:][:lower:] ]*\.
\end{lstlisting}
\end{latin}

این عبارت از سه بخش کلیدی تشکیل شده است:

\begin{itemize}
  \item \cmd{\^{}[[:upper:]]}: لنگر ابتدای خط به همراه کلاس نویسه‌ای برای تضمین شروع با حرف بزرگ.
  \item \cmd{[[:upper:][:lower:] ]*}: یک عبارت براکتی ترکیبی شامل حروف بزرگ، کوچک و نویسه‌ی فاصله که با کمّی‌کننده \cmd{*} همراه شده است تا تکرار این مجموعه را به هر تعداد دفعات (حتی صفر بار) مجاز کند.
  \item \cmd{\textbackslash.}: نویسه‌ی نقطه که با بک‌اسلش گریز داده شده تا از حالت متاکاراکتر خارج و به عنوان یک نقطه‌ی واقعی در انتهای جمله تلقی شود.
\end{itemize}

بیایید این الگو را آزمایش کنیم:

\begin{latin}
\begin{lstlisting}
[me@linuxbox ~]$ echo "This works." | grep -E '^[[:upper:]][[:upper:][:lower:]]*\.'
This works.
[me@linuxbox ~]$ echo "This Works." | grep -E '^[[:upper:]][[:upper:][:lower:]]*\.'
This Works.
[me@linuxbox ~]$ echo "this does not" | grep -E '^[[:upper:]][[:upper:][:lower:]]*\.'
[me@linuxbox ~]$
\end{lstlisting}
\end{latin}

همان‌طور که مشاهده می‌شود، دو مورد نخست به دلیل رعایت ساختار (شروع با حرف بزرگ و پایان با نقطه) تطبیق داده شدند، اما مورد سوم به دلیل رعایت نکردن این ضوابط از خروجی حذف گردید.

\subsection{نماد + (تطبیق یک یا چندباره)}

نماد \cmd{+} عملکردی مشابه با ستاره (\cmd{*}) دارد، با این تفاوت که تکرار عنصر پیشین را حداقل یک بار الزامی می‌کند؛ به عبارت دیگر، وجود آن عنصر در الگو اختیاری نیست.

در ادامه، یک عبارت باقاعده را بررسی می‌کنیم که منحصراً خطوطی را تطبیق می‌دهد که از گروه‌های یک یا چندتایی حروف الفبا تشکیل شده و با «تنها یک فاصله» از هم جدا شده‌اند:

\begin{latin}
\begin{lstlisting}
[me@linuxbox ~]$ echo "This that" | grep -E '^([[:alpha:]]+ ?)+$'
This that
[me@linuxbox ~]$ echo "a b c" | grep -E '^([[:alpha:]]+ ?)+$'
a b c
[me@linuxbox ~]$ echo "a b 9" | grep -E '^([[:alpha:]]+ ?)+$'
[me@linuxbox ~]$ echo "abc  d" | grep -E '^([[:alpha:]]+ ?)+$'
[me@linuxbox ~]$
\end{lstlisting}
\end{latin}

تحلیل رفتار این الگو نتایج جالبی را نشان می‌دهد:

\begin{itemize}
  \item رشته‌ی \cmd{"a b 9"} تطبیق داده نمی‌شود، زیرا حاوی نویسه‌ی عددی است که در کلاس \cmd{[:alpha:]} تعریف نشده است.
  \item رشته‌ی \cmd{"abc  d"} نیز رد می‌شود، زیرا وجود دو فاصله متوالی با بخش \cmd{?} (که حداکثر یک فاصله را مجاز می‌داند) در تضاد است.
\end{itemize}

\subsection{نمادهای \{ و \} (تکرارکننده‌های بازه‌ای یا Bound Quantifiers)}

متاکاراکترهای \cmd{\{} و \cmd{\}} برای تعیین دقیق حداقل و حداکثر تعداد دفعات تکرار یک عنصر در الگوهای جست‌وجو به کار می‌روند. این کمّی‌کننده‌ها طبق استاندارد \en{POSIX}، به چهار شیوه‌ی مختلف که در جدول زیر تشریح شده است، قابل استفاده هستند:

\begin{table}[htbp]
\centering
\caption{تعیین تعداد دفعات تطابق در کمّی‌کننده‌های بازه‌ای}
\label{tab:bound-quantifiers}
\begin{tabularx}{\textwidth}{l X}
\toprule
\textbf{نحوه نگارش} & \textbf{مفهوم فنی} \\
\midrule
\cmd{\{n\}} & تطبیق عنصر پیشین در صورتی که دقیقاً \en{n} بار تکرار شده باشد. \\
\addlinespace
\cmd{\{n,m\}} & تطبیق عنصر پیشین در صورتی که حداقل \en{n} بار و حداکثر \en{m} بار تکرار شده باشد. \\
\addlinespace
\cmd{\{n,\}} & تطبیق عنصر پیشین در صورتی که \en{n} بار یا بیشتر تکرار شده باشد. \\
\addlinespace
\cmd{\{,m\}} & تطبیق عنصر پیشین در صورتی که حداکثر \en{m} بار تکرار شده باشد. \\
\bottomrule
\end{tabularx}
\end{table}

با بازگشت به مثال اعتبارسنجی شماره تلفن، می‌توانیم با بهره‌گیری از این قابلیت، عبارت باقاعده پیچیده‌ی قبلی را:

\begin{latin}
\begin{lstlisting}
^\(?[0-9][0-9][0-9]\)? [0-9][0-9][0-9]-[0-9][0-9][0-9][0-9]$
\end{lstlisting}
\end{latin}

به نسخه‌ای خواناتر و بهینه‌تر تبدیل کنیم:

\begin{latin}
\begin{lstlisting}
^\(?[0-9]{3}\)? [0-9]{3}-[0-9]{4}$
\end{lstlisting}
\end{latin}

اکنون کارکرد این الگوی بازنویسی‌شده را بررسی می‌کنیم:

\begin{latin}
\begin{lstlisting}
[me@linuxbox ~]$ echo "(555) 123-4567" | grep -E '^\(?[0-9]{3}\)? [0-9]{3}-[0-9]{4}$'
(555) 123-4567
[me@linuxbox ~]$ echo "555 123-4567" | grep -E '^\(?[0-9]{3}\)? [0-9]{3}-[0-9]{4}$'
555 123-4567
[me@linuxbox ~]$ echo "5555 123-4567" | grep -E '^\(?[0-9]{3}\)? [0-9]{3}-[0-9]{4}$'
[me@linuxbox ~]$
\end{lstlisting}
\end{latin}

همان‌طور که در خروجی مشاهده می‌شود، عبارت باقاعده جدید به درستی شماره‌هایی را که با ساختار کد ناحیه (با پرانتز یا بدون آن) مطابقت دارند تایید کرده و مواردی را که تعداد ارقام آن‌ها نامعتبر است (مانند پیش‌شماره چهار رقمی) رد می‌کند.

\section{کاربردهای عملی عبارات باقاعده}

در این بخش، برخی از دستوراتی را که پیش‌تر آموخته‌ایم بازبینی کرده و نحوه‌ی ادغام آن‌ها با عبارات باقاعده را بررسی می‌کنیم.

\subsection{اعتبارسنجی فهرست شماره تلفن با ابزار grep}

در مثال‌های پیشین، اعتبارسنجی شماره‌ها را به صورت انفرادی انجام دادیم؛ اما در سناریوهای واقعی، معمولاً با فهرستی از داده‌ها سر و کار داریم. ابتدا با استفاده از یک دستور مرکب در خط فرمان، فایلی حاوی شماره‌های تصادفی ایجاد می‌کنیم. اگرچه برخی از اجزای این دستور در فصل‌های آتی تشریح خواهند شد، اما فعلاً آن را به عنوان یک ابزار تولید داده در نظر بگیرید:

\begin{latin}
\begin{lstlisting}
[me@linuxbox ~]$ for i in {1..10}; do echo "(${RANDOM:0:3}) ${RANDOM:0:3}-${RANDOM:0:4}" >> phonelist.txt; done
\end{lstlisting}
\end{latin}

این دستور فایلی به نام \file{phonelist.txt} ایجاد کرده و ده شماره تلفن با ساختار تصادفی در آن درج می‌کند. با بررسی محتوای فایل، متوجه می‌شویم که برخی از داده‌ها از نظر تعداد ارقام دچار نقص هستند:

\begin{latin}
\begin{lstlisting}
[me@linuxbox ~]$ cat phonelist.txt
(232) 298-2265
(624) 381-1078
(540) 126-1980
(874) 163-2885
(286) 254-2860
(292) 108-518
(129) 44-1379
(458) 273-1642
(686) 299-8268
(198) 307-2440
\end{lstlisting}
\end{latin}

این نقص در داده‌ها، فرصت مناسبی را فراهم می‌کند تا از \cmd{grep} برای شناسایی رکوردهای نامعتبر استفاده کنیم. یک راهبرد مؤثر در اعتبارسنجی، استخراج و نمایش خطوطی است که با استاندارد ما مطابقت ندارند:

\begin{latin}
\begin{lstlisting}
[me@linuxbox ~]$ grep -Ev '^\([0-9]{3}\) [0-9]{3}-[0-9]{4}$'
phonelist.txt
(292) 108-518
(129) 44-1379
[me@linuxbox ~]$
\end{lstlisting}
\end{latin}

در این دستور، از گزینه \cmd{-v} (\en{Invert Match}) استفاده شده است تا \cmd{grep} منحصراً خطوطی را نمایش دهد که با الگوی تعیین‌شده مطابقت ندارند. در عبارت باقاعده فوق، لنگرهای ابتدای خط (\cmd{\^{}}) و انتهای خط (\cmd{\$}) تضمین می‌کنند که کل سطر دقیقاً با الگو منطبق باشد و هیچ نویسه‌ی اضافه‌ای در پیرامون شماره وجود نداشته باشد. برخلاف مثال‌های قبلی، در اینجا پرانتزها به عنوان بخشی الزامی از ساختار شماره در نظر گرفته شده‌اند.

\subsection{شناسایی نام‌فایل‌های غیرمتعارف با ابزار find}

دستور \cmd{find} از معیاری مبتنی بر عبارات باقاعده (\en{Regular Expressions}) برای جست‌وجوی فایل‌ها پشتیبانی می‌کند. با این حال، هنگام به‌کارگیری عبارات باقاعده در \cmd{find} نسبت به ابزار \cmd{grep} یک تفاوت ساختاری کلیدی وجود دارد که باید به آن توجه داشت: در حالی که \cmd{grep} هر خطی را که صرفاً «شامل» الگوی مورد نظر باشد چاپ می‌کند، دستور \cmd{find} الزاماً به دنبال تطبیق «تمامیتِ نام مسیر» (\en{Pathname}) با عبارت باقاعده است.

در مثال زیر، از \cmd{find} به همراه یک عبارت باقاعده برای شناسایی مسیرهایی استفاده می‌کنیم که حاوی نویسه‌هایی خارج از مجموعه‌ی استاندارد زیر باشند:

\begin{latin}
\begin{lstlisting}
[-_./0-9a-zA-Z]
\end{lstlisting}
\end{latin}

اجرای این دستور، فایل‌هایی را که در نام یا مسیر خود دارای فاصله (\en{Space}) یا سایر نویسه‌های بالقوه مخاطره‌آمیز هستند، آشکار می‌کند:

\begin{latin}
\begin{lstlisting}
[me@linuxbox ~]$ find . -regex '.*[^-_./0-9a-zA-Z].*'
\end{lstlisting}
\end{latin}

به دلیل ضرورت تطبیق الگو با کلِ مسیر فایل، از عبارت \cmd{.*} در ابتدا و انتهای الگو استفاده شده است تا هر تعداد نویسه در پیرامون بخش مورد نظر پوشش داده شود. در هسته‌ی این عبارت نیز از یک «مجموعه‌ی نفی‌شده» (\en{Negated Character Set}) استفاده کرده‌ایم؛ این مجموعه تمامی نویسه‌های مجاز و استاندارد را شامل می‌شود، بنابراین هر نویسه‌ای که در این فهرست نباشد، منجر به تطبیق و نمایش مسیر فایل خواهد شد.

\subsection{جست‌وجوی فایل‌ها با ابزار locate}

برنامه‌ی \cmd{locate} از هر دو نوع عبارات باقاعده پایه (با استفاده از گزینه‌ی \cmd{--regexp}) و عبارات باقاعده توسعه‌یافته (با استفاده از گزینه‌ی \cmd{--regex}) پشتیبانی می‌کند. با بهره‌گیری از این قابلیت، می‌توانیم بسیاری از عملیات جست‌وجویی را که پیش‌تر بر روی فایل‌های نمونه‌ی \cmd{dirlist} انجام دادیم، در سطح کل سیستم پیاده‌سازی کنیم.

در مثال زیر، با استفاده از قابلیت تناوب (\en{Alternation})، مسیرهایی را جست‌وجو می‌کنیم که شامل یکی از الگوهای \cmd{bin/bz}، \cmd{bin/gz} یا \cmd{bin/zip} باشند:

\begin{latin}
\begin{lstlisting}
[me@linuxbox ~]$ locate --regex 'bin/(bz|gz|zip)'
/bin/bzcat
/bin/bzcmp
/bin/bzdiff
/bin/bzegrep
/bin/bzexe
/bin/bzfgrep
/bin/bzgrep
/bin/bzip2
/bin/bzip2recover
/bin/bzless
/bin/bzmore
/bin/gzexe
/bin/gzip
/usr/bin/zip
/usr/bin/zipcloak
/usr/bin/zipgrep
/usr/bin/zipinfo
/usr/bin/zipnote
/usr/bin/zipsplit
\end{lstlisting}
\end{latin}

در این الگو، استفاده از پرانتز باعث می‌شود که عملگر تناوب (\cmd{|}) منحصراً بر روی پیشوندهای مشخص‌شده اعمال شود. این روش ابزاری بسیار سریع و کارآمد برای مکان‌یابی دسته‌جمعی فایل‌های مرتبط با ابزارهای فشرده‌سازی در دایرکتوری‌های سیستمی است.

\subsection{جست‌وجوی متنی در محیط‌های less و vim}

ابزارهای \cmd{less} و \cmd{vim} از مکانیزم مشابهی برای پیمایش و جست‌وجوی متون استفاده می‌کنند. در هر دو برنامه، با فشردن کلید \cmd{/} و وارد کردن یک عبارت باقاعده (\en{Regular Expression})، فرآیند تطبیق الگو آغاز می‌شود. اگر فایل \file{phonelist.txt} را جهت بازبینی در \cmd{less} فراخوانی کنیم:

\begin{latin}
\begin{lstlisting}
[me@linuxbox ~]$ less phonelist.txt
\end{lstlisting}
\end{latin}

سپس الگوی اعتبارسنجی خود را برای یافتن شماره‌های استاندارد وارد نماییم:

\begin{latin}
\begin{lstlisting}
(232) 298-2265
(624) 381-1078
(540) 126-1980
(874) 163-2885
(286) 254-2860
(292) 108-518
(129) 44-1379
(458) 273-1642
(686) 299-8268
(198) 307-2440
~
~
~
/^\([0-9]{3}\) [0-9]{3}-[0-9]{4}$
\end{lstlisting}
\end{latin}

ابزار \cmd{less} رشته‌های منطبق با الگو را برجسته (\en{Highlight}) می‌کند. این قابلیت باعث می‌شود تا رکوردهای فاقد قالب‌بندی صحیح که برجسته نشده‌اند، به سرعت شناسایی شوند.

\begin{latin}
\begin{lstlisting}
(232) 298-2265
(624) 381-1078
(540) 126-1980
(874) 163-2885
(286) 254-2860
(292) 108-518
(129) 44-1379
(458) 273-1642
(686) 299-8268
(198) 307-2440
~
~
~
(END)
\end{lstlisting}
\end{latin}

در مقابل، ویرایشگر \cmd{vim} به‌صورت پیش‌فرض از عبارات باقاعده پایه (\en{BRE}) استفاده می‌کند؛ بنابراین عبارت جست‌وجوی ما در این محیط به شکل زیر تغییر می‌یابد:

\begin{latin}
\begin{lstlisting}
/\([0-9]\{3\}\) [0-9]\{3\}-[0-9]\{4\}
\end{lstlisting}
\end{latin}

همان‌طور که مشاهده می‌شود، منطق الگو ثابت مانده است، اما بر اساس قواعد \en{BRE}، بسیاری از نویسه‌هایی که در حالت توسعه‌یافته (\en{ERE}) به‌عنوان متاکاراکتر شناخته می‌شوند، در اینجا نویسه‌ی معمولی قلمداد شده و تنها در صورتی که با بک‌اسلش (\cmd{\textbackslash}) گریز داده شوند، معنای فنی و عملیاتی پیدا می‌کنند.

بسته به پیکربندی اختصاصی ویرایشگر \cmd{vim} در سیستم شما، نتایج منطبق با الگو به‌صورت خودکار برجسته (\en{Highlight}) می‌شوند. در صورتی که این ویژگی فعال نیست، می‌توانید با اجرای دستور زیر در حالت خط فرمان (\en{Command Mode})، برجسته‌سازی را فعال نمایید:

\begin{latin}
\begin{lstlisting}
:hlsearch
\end{lstlisting}
\end{latin}

\begin{note}
\textbf{نکته:} قابلیت برجسته‌سازی نتایج جست‌وجو ممکن است بسته به توزیع لینوکس مورد استفاده و نسخه‌ی نصب‌شده‌ی \cmd{vim} متفاوت باشد. برای مثال، توزیع اوبونتو به‌صورت پیش‌فرض نسخه‌ای محدود و سبک از این ویرایشگر را ارائه می‌دهد که فاقد برخی ویژگی‌های پیشرفته است. در چنین شرایطی، توصیه می‌شود با استفاده از مدیر بسته‌ی سیستم خود (\en{Package Manager})، نسخه‌ی کامل‌تر را نصب کنید تا از تمامی قابلیت‌های عملیاتی آن بهره‌مند شوید.
\end{note}

\section{جمع‌بندی}

در این فصل، با بخشی از کاربردهای گسترده و متنوع عبارات باقاعده (\en{Regular Expressions}) آشنا شدیم. برای کشف ابزارهای بیشتری که از این قابلیت پشتیبانی می‌کنند، می‌توانیم با استفاده از خودِ عبارات باقاعده در صفحات راهنما (\en{man pages}) جست‌وجو کنیم:

\begin{latin}
\begin{lstlisting}
[me@linuxbox ~]$ cd /usr/share/man/man1
[me@linuxbox man1]$ zgrep -El 'regex|regular expression' *.gz
\end{lstlisting}
\end{latin}

ابزار \cmd{zgrep} به‌عنوان یک واسط (\en{Frontend}) برای \cmd{grep} عمل کرده و امکان جست‌وجوی مستقیم در فایل‌های فشرده را فراهم می‌کند. در دستور فوق، ما بخش \en{1} صفحات راهنما را در مسیر استاندارد آن‌ها پیمایش کردیم. خروجی این دستور، فهرستی از برنامه‌هایی است که در مستندات خود به مفاهیم \en{regex} یا \en{regular expression} اشاره کرده‌اند. تکرار این مفاهیم در مستندات سیستمی، گواهی بر اهمیت و نفوذ بالای عبارات باقاعده در ابزارهای مختلف لینوکس است.

در فصل آتی، به بررسی یکی از ویژگی‌های پیشرفته در عبارات باقاعده پایه (\en{BRE}) به نام ارجاع‌های بازگشتی (\en{Backreferences}) خواهیم پرداخت که امکان بازخوانی و انطباق مجدد الگوهای یافت‌شده را فراهم می‌سازد.

\section{منابع مطالعه تکمیلی}

برای تسلط بیشتر بر عبارات باقاعده، منابع آنلاین متعددی در دسترس هستند. علاوه بر این، مطالعه مقالات زیر در ویکی‌پدیا برای درک عمیق‌تر مفاهیم بنیادین مرتبط توصیه می‌شود:

\en{POSIX}:

\begin{latin}
\url{http://en.wikipedia.org/wiki/Posix}
\end{latin}

\en{ASCII}:

\begin{latin}
\url{http://en.wikipedia.org/wiki/Ascii}
\end{latin}